\chapter{Gen5.1 Native：实现解剖与代码印证}
\label{ch:native-impl}

\section{代码布局}

\begin{longtable}{@{}p{5.6cm}p{8.1cm}@{}}
\toprule
路径 & 职责 \\
\midrule
\endhead
\filepath{topo\_phn.h} & 公开配置 / Slice \\
\filepath{topo\_phn\_internal.h} & Embed、Flip、PH、标定、resume \\
\filepath{topo\_phn\_scan.cc} & 热路径、闭合 Embed、分层 UF \\
\filepath{topo\_phn\_streaming.cc} & 512KB parity \\
\filepath{topo\_phn\_config.cc} & profile / 超块 \\
\filepath{backup\_pipeline.cc} & 三档 runtime 缓存 \\
\bottomrule
\end{longtable}

\section{公开配置}

\begin{lstlisting}[caption={TopoPhnConfig},label={lst:phn-cfg}]
enum class TopoPhnKernel : uint8_t {
  kTriNative = 5, kPhH0Native = 6
};
struct TopoPhnConfig {
  uint32_t min_size, avg_size, max_size;
  TopoPhnKernel kernel{TopoPhnKernel::kTriNative};
  uint32_t table_seed{0};
  uint32_t event_stride{64};  // pow2 after RoundEventStridePow2
  uint8_t  k_points{16};
  uint32_t q_mod{1024};
  uint32_t flip_tau{1};       // runtime re-derive; not in SB
  uint8_t  persist_delta{1};
  bool enable_persist_delta{true};
};
\end{lstlisting}

\section{\texorpdfstring{LandmarkEmbedY4：闭合 = 迭代}{LandmarkEmbedY4 闭合等于迭代}}

\begin{lstlisting}[caption={bit-identical 闭合形式（\filepath{topo\_phn\_scan.cc}）},label={lst:embed-y4}]
constexpr uint32_t kLcgM = 1664525u;
constexpr uint32_t kLcgC = 1013904223u;
inline uint32_t LandmarkEmbedY4(uint32_t seed, uint8_t b0, uint8_t b1,
                                uint8_t b2, uint8_t b3) {
  uint32_t y = seed * kLcgM4 + kLcgCGeom;
  y += uint32_t(b0) * kLcgM3;
  y += uint32_t(b1) * kLcgM2;
  y += uint32_t(b2) * kLcgM;
  y += uint32_t(b3);
  y ^= uint32_t(b2) << 8;
  y ^= uint32_t(b1) << 16;
  y ^= uint32_t(b0) << 24;
  return y;
}
uint32_t LandmarkEmbedY(...) {
  if (p >= 3)
    return LandmarkEmbedY4(seed, data[p-3], data[p-2], data[p-1], data[p]);
  // short prefix: iterative LCG (rare)
}
\end{lstlisting}

单测强制闭合形式与 4 步迭代一致；AVX2 据此向量化邻域加载。

\section{pow2 stride}

\begin{lstlisting}[caption={最近二的幂（平局偏小）},label={lst:pow2-stride}]
inline uint32_t RoundEventStridePow2(uint32_t s) {
  s = ClampEventStride(s); // [8, 2MiB]
  // find hi = next pow2 >= s, lo = hi/2
  // return (s-lo <= hi-s) ? lo : hi;
}
// Soft floor for PH densify:
inline uint32_t PhnProductiveStrideFloor(uint32_t avg_size) {
  return RoundEventStridePow2(max(32u, avg_size / 1024u));
}
\end{lstlisting}

扫描：$Y\mathbin{\&}(S-1)=0$ 替代昂贵取模。

\section{FlipCountFixed 完整逻辑}

\begin{lstlisting}[caption={Native Flip（环序 + CrossI）},label={lst:phn-flip}]
inline int64_t CrossI(int64_t ax, int64_t ay, int64_t bx, int64_t by) {
  return ax * by - ay * bx;
}
uint32_t FlipCountFixed(..., uint8_t ring_start) {
  // Y(i), T(i) via LandmarkIndex(..., ring_start)
  // fan from P0; then consecutive triples; Cross <= 0 => ++flips
  // coords: T % q_mod, Y % q_mod
}
\end{lstlisting}

\begin{figure}[htbp]
\centering
\begin{tikzpicture}[scale=0.9]
  \coordinate (A) at (0,0);
  \coordinate (B) at (1.8,1.2);
  \coordinate (C) at (2.6,0.2);
  \coordinate (D) at (1.0,-0.8);
  \foreach \p/\lab in {A/P_0,B/P_i,C/P_{i+1},D/P_j} {
    \fill (\p) circle (1.5pt);
    \node[font=\scriptsize, above right] at (\p) {$\lab$};
  }
  \draw[-{Stealth}] (A)--(B);
  \draw[-{Stealth}] (A)--(C);
  \node[font=\footnotesize, gray] at (3.8,0.5) {$\mathrm{Cross}\le0\Rightarrow$ flip};
\end{tikzpicture}
\caption{锚点扇形方向谓词（示意）。}
\label{fig:cross-orient}
\end{figure}

\section{分层 PH-H0}

\begin{lstlisting}[caption={嵌套 eps 一次扫边},label={lst:phn-h0}]
constexpr uint32_t kEpsR2[3] = {1u, 4u, 16u};
// Build all pairwise d2; sort ascending
// UF: comps = k; for each eps level, unite while d2 <= R^2
// snapshot out_beta0[ei] = comps
void PhH0BettiVec(...);
bool PhH0Delta(const uint32_t* prev, const uint32_t* cur);
uint8_t PhH0PersistSpan(const uint32_t* beta0); // max-min
\end{lstlisting}

\section{事件评估与切后复位}

\begin{lstlisting}[caption={EvalEvent + ResetLandmarkWindow},label={lst:eval-reset}]
// Tri: push -> |flip_after - prev| >= tau
// PH : beta_before, push, beta_after;
//      flipOK && PhH0Delta [&& PersistSpan >= delta]
bool EvalEventAtLandmark(...);

// After cut: clear ring so next chunk is content-local
void ResetLandmarkWindow(TopoPhnResume* resume) {
  recent_count = 0; ring_start = 0; has_prev = false;
  prev_flip = 0; memset(prev_beta0, 0, ...);
}
\end{lstlisting}

\section{热路径结构}

\begin{figure}[htbp]
\centering
\begin{tikzpicture}[
  box/.style={rectangle, rounded corners=2pt, draw=Gen51Color, fill=Gen51Color!10,
    align=center, font=\footnotesize, minimum height=0.8cm},
  arr/.style={-{Stealth}, thick}]
  \node[box] (a) {bytes};
  \node[box, right=0.5cm of a] (b) {EmbedY4};
  \node[box, right=0.5cm of b] (c) {pow2\&};
  \node[box, right=0.5cm of c] (d) {LM?};
  \node[box, right=0.5cm of d] (e) {Flip/PH};
  \node[box, right=0.5cm of e] (f) {cut+reset};
  \draw[arr] (a)--(b)--(c)--(d)--(e)--(f);
  \draw[arr,dashed] (d.south)--++(0,-0.5)-| node[pos=0.2,below,font=\tiny]{N} (b.south);
\end{tikzpicture}
\caption{绝大多数字节仅 Embed+mask。}
\label{fig:phn-hotpath}
\end{figure}

加速要点（切点不变）：
\begin{enumerate}[nosep]
  \item AVX2：16B 双 8-lane；无命中跳过 store；
  \item Embed 用 \texttt{loadl}+\texttt{cvtepu8}，无 gather；
  \item Pipeline 三档缓存 \texttt{CalibratePhnRuntimeParams} 结果。
\end{enumerate}

\section{Resume 状态}

\begin{lstlisting}[caption={TopoPhnResume},label={lst:phn-resume}]
struct TopoPhnResume {
  size_t scan_rel; bool in_scan;
  uint32_t recent_y[16], recent_t[16];
  uint32_t recent_count; uint8_t ring_start;
  uint32_t prev_flip; uint32_t prev_beta0[3];
  bool has_prev;
  size_t last_cut_abs; bool has_last_cut;
  // rolling_mix retained for ABI; embed is content-local
};
\end{lstlisting}

\texttt{MinGapOk}：距上次切 $\ge\min$。多 feed 与整文件 bit-identical。

\section{测试印证清单}

\begin{itemize}[nosep]
  \item Embed 闭合 $=$ 迭代；
  \item Tri$\neq$PH 切点集；
  \item 流式多 feed parity；
  \item \texttt{reuse\_1byte}；
  \item \texttt{TopoPhnTriTest.*} / \texttt{TopoPhnH0Test.*}。
\end{itemize}
